\documentclass[12pt]{article}
%---DOCUMENT MARGINS---
\usepackage{geometry} % Required for adjusting page dimensions and margins
\geometry{
	paper=a4paper, % Paper size, change to letterpaper for US letter size
	top=4cm, % Top margin
	bottom=2cm, % Bottom margin
	left=2cm, % Left margin
	right=2cm, % Right margin
	headheight=3cm, % Header height
	%footskip=1.5cm, % Space from the bottom margin to the baseline of the footer
	headsep=0.5cm, % Space from the top margin to the baseline of the header
	%showframe, % Uncomment to show how the type block is set on the page
}
\usepackage{cmap}

\usepackage[T1,T2A]{fontenc}
\usepackage{fontspec}
\setmainfont[
	BoldFont={* Bold},
	ItalicFont={* Italic},
	BoldItalicFont={* BoldItalic}
]{DejaVu Serif Condensed}
\setsansfont{DejaVu Sans}
\setmonofont{Dejavu Sans Mono}
\usepackage[math-style=TeX]{unicode-math}
\setmathfont{Latin Modern Math}

\usepackage[nobottomtitles*]{titlesec}
\titleformat
{\section} % command
{\fontsize{14}{14}\sffamily\bfseries} % format
{} % label
{0pt} % sep
{} % before-code
\titlespacing{\section}{0pt}{0.75em}{0.25em}
\newcommand{\tl}{}
\newcommand{\ml}{}
\newcommand{\problem}[1]{%
	\section[#1]{#1 \fontsize{14}{14}\selectfont{\hfill{\normalfont{
				\emoji{hourglass-not-done}\tl \space\space \emoji{floppy-disk} \ml}}}}
}
\AddToHook{cmd/section/before}{\clearpage}

\usepackage[dvipsnames]{xcolor}
\titleformat
{\subsection} % command
{\fontsize{14}{14}\sffamily\bfseries} % format
{\includesvg[width=0.035\textwidth, inkscapelatex=false]{lion}\space} % label
{0pt} % sep
{} % before-code
\titlespacing{\subsection}{0pt}{0.75em}{0.25em}

\setlength{\parskip}{0.5em}
\setlength{\parindent}{0pt}
\renewcommand{\baselinestretch}{1.2}
\sloppy

\usepackage{listings}
\definecolor{lstbg}{RGB}{250,250,252}
\definecolor{lstframe}{RGB}{220,220,225}
\lstdefinestyle{mystyle}{
	backgroundcolor=\color{lstbg},
	frame=single,
	rulecolor=\color{lstframe},
	basicstyle=\ttfamily\small,
	keywordstyle=\bfseries,
	breaklines=true,
	aboveskip=0.5\baselineskip,
	belowskip=0\baselineskip  
}
\lstset{style=mystyle, language=C++}

\usepackage{fancyhdr}
\pagestyle{fancy}
\setlength{\headwidth}{\textwidth}
\newcommand{\round}{}
\newcommand{\problemName}{}
\newcommand{\lang}{English}
\fancyhead[L]{
	\begin{minipage}{0.4\textwidth}
		\bf{
			EJOI 2025 \round\\
			\problemName\\
			\emoji{globe-with-meridians} \lang
		}
	\end{minipage}
	\vspace{0.5cm}
}
\fancyhead[R]{%
	\begin{minipage}{0.6\textwidth}
		\includesvg[width=\textwidth, inkscapelatex=false]{logo}
	\end{minipage}
	\vspace{0.5cm}
}
\usepackage{lastpage}
\fancyfoot[C]{\problemName\space(\lang) \thepage\ / \pageref*{LastPage}}

\raggedbottom

\usepackage{amsmath}
\usepackage{stmaryrd}

\usepackage{graphicx}
\graphicspath{{./}}
\usepackage[export]{adjustbox}
\usepackage{wrapfig}
\makeatletter
\patchcmd\WF@putfigmaybe{\lower\intextsep}{}{}{\fail}
\AddToHook{env/wrapfigure/begin}{\setlength{\intextsep}{0pt}}
\makeatother
\usepackage[inkscapearea=page, inkscapepath=./svg-inkscape]{svg}
\svgpath{{./}}

\usepackage{makecell}
\usepackage{tabularray}
\AtBeginEnvironment{table}{\vspace{-0.2cm}}
\AtEndEnvironment{table}{\vspace{-0.2cm}}
\usepackage{float}
\usepackage{placeins}
\usepackage{caption}
\captionsetup[table]{
	skip=0.25em,
	singlelinecheck=false, justification=justified
}

\usepackage{enumitem}
\setlist{itemsep=-0.4em, leftmargin=22pt, topsep=-\parskip}
\AfterEndEnvironment{itemize}{\vspace{\parskip}}
\newcommand{\tabitem}{\indent~~\llap{\textbullet}~~}

\usepackage{hyperref}
\hypersetup{
	colorlinks=true,
	citecolor=blue,
	linkcolor=blue,
	urlcolor=cyan,
}

\usepackage{emoji}

\usepackage[english]{babel}

\begin{document}

\renewcommand{\round}{Day 2}
\renewcommand{\problemName}{Task Navigation}
\renewcommand{\lang}{Ukrainian}

\renewcommand{\tl}{$3$ sec.}
\renewcommand{\ml}{$512$ MB}
\problem{Задача Навігація}

Дано \textbf{зв'язний неорієнтований простий граф-кактус}\textsuperscript{1} з $N \leq 1000$ вершинами та $M$ ребрами.
Його вершини мають кольори (позначені невід'ємними числами від $0$ до $1499$).
Спочатку всі вершини мають колір 0.
\textbf{Детерменістичний робот без пам'яті}\textsuperscript{2} досліджує граф ходячи від вершини до вершини.
Він має відвідати всі вершини хоча б один раз і потім зупинитися.

Робот починає в якійсь вершині, це може бути будь-яка з вершин графу.
На кожному кроці він бачить колір вершини в якій знаходиться та кольори всіх сусідніх вершин
\textbf{в порядку, який для теперешньої вершини є фіксованим}
(тобто при повторному відвідуванні цієї вершини робот отримає таку саму послідовність сусідніх вершин,
навіть якщо їх кольори змінилися).
Робот робить одну з наступних дій:

\begin{enumerate}
    \item Вирішує зупинитися.
    \item Вибирає колір для поточної вершини(можливо такий самий як зараз) і в яку сусідню вершину перейти. 
    Та сусідня вершина ідентифікується індексом від $0$ до $D-1$, де $D$ це кількість сусідніх вершин.
\end{enumerate}

В другому випадку поточна вершина перефарбовується (або можливо не  перефарбовується)
і робот переходить в обрану сусідню вершину. Це повторюється до того часу, поки робот не вирішить зупинитися або не привищить ліміт кроків. Робот виграє якщо він відвідує всі всі вершини а потім зупиняється в межах ліміту $L = 3000$ кроків (інакше він програє).


Ви повинні розробити стратегію для робота, яка дозволяє вирішити цю задачу на будь-якому какстусі, який задовольняє обмеженням. Також, ви повинні спробувати мінімізувати кількість різних кольорів, які використовуються в вашому рішенні. Колір 0 завжди враховується як використаний.

\textsuperscript{1}\textit{Зв'язний неорієнтований простий граф-кактус це зв'язний неорієнтований простий (Кожна вершина досяжна; ребра двонаправленні; в графі відсутні петлі та мульти-ребра) в якому кожне ребро належить не більш ніж одному простому циклу. (Простий цикл -- цикл в якому кожна вершина зустрічається не більше одного разу). Знизу є приклад.}

\begin{adjustbox}{center}
	\includesvg[inkscapelatex=false, width=.5\linewidth]{graph}
\end{adjustbox}

\textsuperscript{2}\textit{Робот називається детерміністичним та без пам'яті, якщо його дії залежать тільки від його вхідних даних (тобто він не зберігає дані між кроками) і він завжди виконує одну і ту ж саму дію, коли отримає однакові вхідні дані.}

\subsection{Деталі реалізації}
Стратегія робота повинна бути реалізована як функція:
\begin{lstlisting}
 std::pair<int, int> navigate(int currColor, std::vector<int> adjColors)
\end{lstlisting}

Вона отримує колір поточної вершини та кольори всіх її сусідів (в порядку) як аргументи. Вона має повертати пару, перший елемент якої це новий колір тепершньої вершини, а другий -- індекс в списку сусідів вершини, в яку робот має перейти. Якщо робот має зупинитися, то ця функція повинна повернути $(-1, -1)$

Ця функція буде викликатися неодноразово, щоб вибрати дії робота. Оскільки вона детермінована, якщо \texttt{navigate} вже була викликана з деякими аргументами, вона ніколи більше не буде викликана з тими самими параметрами; натомість буде повторно використано її попередньо повернене значення. Крім того, кожен тест може містити $T \leq 5$ підтестів (окремих графів та/або початкових позицій), і вони можуть виконуватися одночасно (тобто ваша програма може отримувати чергуючі виклики з різних підтестів). Нарешті, виклики \texttt{navigate} можуть відбуватися в \textbf{окремих виконаннях} вашої програми (але іноді вони також можуть відбуватися в одному виконанні). Загальна кількість виконань вашої програми становить $P = 100$. Через усе це ваша програма не повинна пробувати передавати інформацію між різними викликами.

\subsection{Обмеження}

\begin{itemize}
\item $3 \leq N \leq 1000$
\item $0 \leq \mathrm{Color} < 1500$
\item $L = 3000$
\item $T \leq 5$
\item $P = 100$
\end{itemize}

\subsection{Оцінювання}

Частка балів $S$ від макисальної кількості балів за підзадачу залежить від $C$ -- максимальної кількості різних кольорів які ваше рішення використовує (включаючи колір $0$) серед тестів цієї підзадачі і всіх необхідних для неї підзадач:

\begin{itemize}
\item Якщо ваше рішення не працюяє на якомусь з тестів, то $S = 0$.
\item Якщо $C \leq 4$, то $S = 1.0$.
\item Якщо $4 < C \leq 8$, то $S = 1.0 - 0.6 \frac{C - 4}{4}$.
\item If $8 < C \leq 21$, то $S = 0.4 \frac{8}{C}$.
\item If $C > 21$, то $S = 0.15$.
\end{itemize}

\subsection{Підзадачі}

\begin{table}[H]
\begin{tblr}{|Q[c,m]|Q[c,m]|Q[c,m]|Q[c,m]|X[c,m]|}
\hline
\textbf{Підзадача} & \textbf{Бали} & \textbf{\makecell{Необхідні\\підзадачі}} & $N$ & \textbf{Додаткові обмеження}\\
\hline
$0$ & $0$ & $-$ & $\leq 300$ & Приклад. \\
\hline
$1$ & $6$ & $-$ & $\leq 300$ & Граф -- цикл.\textsuperscript{1} \\
\hline
$2$ & $7$ & $-$ & $\leq 300$ & Граф -- зірка.\textsuperscript{2} \\ 
\hline
$3$ & $9$ & $-$ & $\leq 300$ & Граф -- шлях.\textsuperscript{3} \\ 
\hline
$4$ & $16$ & $2 - 3$ & $\leq 300$ & Граф -- дерево.\textsuperscript{4} \\
\hline
$5$ & $27$ & $-$ & $\leq 300$ & Всі вершини мають не більше ніж $3$ сусідні вершини і вершина, в якій робот починає шлях, має $1$-го сусіда. \\
\hline
$6$ & $28$ & $0 - 5$ & $\leq 300$ & $-$ \\
\hline
$7$ & $7$ & $0 - 6$ & $-$ & $-$ \\
\hline
\end{tblr}
\caption*{\textsuperscript{1}Цикл має ребра: $\left(i, (i+1) \bmod N\right)$ для всіх  $0 \leq i < N$. \\
\textsuperscript{2}Зірка має ребра: $\left(0, i\right)$ для всіх  $1 \leq i < N$. \\
\textsuperscript{3}Шлях має ребра: $\left(i, i+1\right)$ для всіх $0 \leq i < N - 1$. \\
\textsuperscript{4}Дерево це граф без циклів.}
\end{table}
\FloatBarrier

\subsection{Приклад}

Розглянемо приклад графу з зображення в умові, який має $N=7$, $M=8$ і ребра $(0, 1)$, $(1, 2)$, $(2, 0)$, $(2, 3)$, $(3, 4)$, $(4, 2)$, $(3, 5)$ та $(2, 6)$. Оскільки порядок елементів в списку сусідів вершини важливий, ми вказали його в цій таблиці:

\begin{table}[H]
\begin{tblr}{|Q[c,m]|Q[c,m]|}
\hline
\textbf{\makecell{Вершина}} & \textbf{Сусідні Вершини} \\
\hline
$0$ & $2, 1$ \\
\hline
$1$ & $2, 0$ \\
\hline
$2$ & $0, 3, 4, 6, 1$ \\
\hline
$3$ & $4, 5, 2$ \\
\hline
$4$ & $2, 3$ \\
\hline
$5$ & $3$ \\
\hline
$6$ & $2$ \\
\hline
\end{tblr}
\end{table}

Припустимо що робот починає в вершині $5$. Тоді могла б бути така (не успішна) послідовність кроків:

\begin{table}[H]
\begin{tblr}{|Q[c,m]|Q[c,m]|Q[c,m]|X[c,m]|Q[c,m]|}
\hline
\textbf{\makecell{\#}} & \textbf{\makecell{Колір}} & \textbf{\makecell{Вершина}} & \textbf{\makecell{Виклик \texttt{navigate}}} & \textbf{Повернене значення} \\
\hline
$1$ & $0, 0, 0, 0, 0, 0, 0$ & $5$ & \texttt{navigate(0, \{0\})} & \texttt{\{1, 0\}} \\
\hline
$2$ & $0, 0, 0, 0, 0, 1, 0$ & $3$ & \texttt{navigate(0, \{0, 1, 0\})} & \texttt{\{4, 2\}} \\
\hline
$3$ & $0, 0, 0, 4, 0, 1, 0$ & $2$ & \texttt{navigate(0, \{0, 4, 0, 0, 0\})} & \texttt{\{0, 3\}} \\
\hline
$4$ & $0, 0, 0, 4, 0, 1, 0$ & $6$ & \textsuperscript{1}\texttt{navigate(0, \{0\})} & \texttt{\{1, 0\}} \\
\hline
$5$ & $0, 0, 0, 4, 0, 1, 1$ & $2$ & \texttt{navigate(0, \{0, 4, 0, 1, 0\})} & \texttt{\{8, 0\}} \\
\hline
$6$ & $0, 0, 8, 4, 0, 1, 1$ & $0$ & \texttt{navigate(0, \{8, 0\})} & \texttt{\{3, 0\}} \\
\hline
$7$ & $3, 0, 8, 4, 0, 1, 1$ & $2$ & \texttt{navigate(8, \{3, 4, 0, 1, 0\})} & \texttt{\{2, 2\}} \\
\hline
$8$ & $3, 0, 2, 4, 0, 1, 1$ & $4$ & \texttt{navigate(0, \{2, 4\})} & \texttt{\{1, 1\}} \\
\hline
$9$ & $3, 0, 2, 4, 1, 1, 1$ & $3$ & \texttt{navigate(4, \{1, 1, 2\})} & \texttt{\{-1, -1\}} \\
\hline
\end{tblr}
\end{table}

Тут робот використав загалом $6$ різних кольорів: $0$, $1$, $2$, $3$, $4$ та $8$ (зауважте, що $0$ вважався б використаним, навіть якби робот ніколи не повертав колір $0$, оскільки всі вершини починаються з кольору $0$). Робот виконав $9$ кроків, перш ніж зупинитися. Однак він зазнав невдачі, оскільки зупинився, не відвідавши вершину $1$.

\textsuperscript{1}Зверніть увагу, що виклик \texttt{navigate} на кроці $4$ насправді не відбудеться. Це тому, що він еквівалентний виклику на кроці $1$, тому градер просто повторно використає повернене значення вашої функції з того виклику. Однак це все одно вважається як крок робота.
\subsection{Приклад градера}

Зразковий градер не робить кілька виконань вашої програми, тому всі виклики \texttt{navigate} будуть в одному виконанні вашої програми.

Формат вхідних даних наступний: Спочатку зчитується $T$ (кількість підтестів). Потім для кожного підтесту:
\begin{itemize}
\item рядок $1$: дві цілих числа -- $N$ та $M$;
\item рядок $2+i$ (для $0 \leq i < M$): два числа -- $A_i$ та $B_i$, які є вершинами з'єднані ребром $i$ ($0 \leq A_i, B_i < N$).
\end{itemize}

Потім зразковий градер виведе кількість різних кольорів, використаних у вашому розв'язку, та кількість кроків, яка була йому потрібна перед зупинкою. Або ж виведе повідомлення про помилку, якщо ваше розв'язання не вдалося.

За замовчуванням, зразковий градер друкує детальну інформацію про те, що бачить і робить робот на кожному кроці. Ви можете вимкнути це, змінивши значення \texttt{DEBUG} з \texttt{true} на \texttt{false}.

\end{document}
